\section{Proofs}\label{app:proofs}
% Translated from G4/PROOFS.md (self-reviewed; independently reviewed 26/10/2026: 0 critical, 4 major, 12 minor, all
% addressed). Keep in sync with that file. Citations marked VERIFY in refs.bib must be checked (G4-6) before submission.

\paragraph{Assumptions.} (A0) $\W\sim p(\W\mid X,y)$, independent of $(x_*,y_*)$ given $(X,y)$; nuisance parameters of
the observation model, if any, are a priori independent of $(\phi,X,y)$. (A1) $\tfm(\cdot\mid x,D)=\ppd(\cdot\mid x,D)$ for
every context the algorithm queries. (A2) $\ppd(k\mid x,D)>0$ for all $k$ and every context of positive probability.
(A3) $\lik$ is known and $\sum_y p(y\mid X)\lik(y)>0$. For instance-level observations
$\lik(y)=\prod_i\lik_i(y_i)$; the set of free points is $F=\{i:\lvert\{k:\lik_i(k)>0\}\rvert\ge2\}$.

\subsection{Lemma 1}
\begin{proof}[Proof of \cref{lem:fullcond}]
Given $\phi$ the pairs $(x_j,y_j)$ are i.i.d., hence exchangeable, so the joint law of $(D_{-i},(x_i,y_i))$ equals that of
$(D_{n-1},(x_*,y_*))$, a context of $n-1$ points and a query. The information $(X,y_{-i})$ coincides with $(x_i,D_{-i})$.
Therefore
\[p(y_i\mid y_{-i},X)=\frac{p((x_i,y_i),D_{-i})}{\sum_kp((x_i,k),D_{-i})}=p(y_i\mid x_i,D_{-i}).\qedhere\]
\end{proof}
The PPD conditions on the query features $x_*$ as well, which may carry information about $\phi$; this is the quantity
the PFN training loss targets, since $x_*$ is drawn from the prior together with the context \citep{muller2022pfn}.

\begin{corollary}[Positivity]\label{cor:positivity}
Under (A2), $p(y\mid X)>0$ for every $y\in[C]^n$.
\end{corollary}
\begin{proof}
Some $y^0$ has $p(y^0\mid X)>0$. Change one coordinate at a time, $y^0\to y^1\to\dots\to y^m=y$, step $k$ changing
point $s_k$. If $p(y^{k-1}\mid X)>0$, then $p(y^{k-1}_{-s_k}\mid X)>0$, the context $D^{k-1}_{-s_k}$ has positive
probability, and by \cref{lem:fullcond} and (A2)
$p(y^k\mid X)=p(y^{k-1}_{-s_k}\mid X)\,p(y^k_{s_k}\mid x_{s_k},D^{k-1}_{-s_k})>0$. Induct on $k$.
\end{proof}

\subsection{Theorem 1}
\begin{proof}[Proof of \cref{thm:exact}(a)]
The update of point $i$ is $P_i(y,y')=\pi_i(y'_i\mid y_{-i})\mathbb 1[y'_{-i}=y_{-i}]$ with
$\pi_i(k\mid y_{-i})\propto\tfm(k\mid x_i,D_{-i})\lik_i(k)$. By \cref{lem:fullcond} and (A1),
$\pi_i(k\mid y_{-i})\propto p(k\mid y_{-i},X)\lik_i(k)\propto\pi(k,y_{-i})$, the full conditional of $\pi$; hence
$\pi(y)P_i(y,y')=\pi(y_{-i})\pi_i(y_i\mid y_{-i})\pi_i(y'_i\mid y_{-i})$ is symmetric in $(y,y')$. By
\cref{cor:positivity}, $\pi(y)>0$ iff $\lik(y)>0$, a product set; any two of its points are joined by changing one
coordinate at a time inside the set, each step with probability at least $\pi_i(\cdot\mid\cdot)/\lvert F\rvert>0$, so the
random-scan kernel $P=\lvert F\rvert^{-1}\sum_{i\in F}P_i$ is irreducible, and aperiodic since $P(y,y)>0$. A finite,
irreducible, aperiodic chain converges geometrically. The implementation sweeps the free points in a uniformly random
order, $P_{\mathrm{sw}}=\frac1{\lvert F\rvert!}\sum_\sigma P_{\sigma(1)}\cdots P_{\sigma(\lvert F\rvert)}$; each factor
preserves $\pi$, and $P_{\mathrm{sw}}$ is reversible because the adjoint of a product is the product in reverse order and
the set of permutations is closed under reversal.
\end{proof}

\begin{proof}[Proof of \cref{thm:exact}(b)]
Choose a bag uniformly; if it is not pure, choose uniformly an ordered pair $(i,j)$ in it with $y_i\neq y_j$ and propose
$y'$ by swapping their labels. Brook's lemma through the intermediate state $(y'_i,y_{-i})$ gives
\[
\frac{p(y'\mid X)}{p(y\mid X)}=\frac{p(y'_i\mid y_{-i},X)}{p(y_i\mid y_{-i},X)}\cdot\frac{p(y'_j\mid y'_i,y_{-ij},X)}{p(y_j\mid y'_i,y_{-ij},X)},
\]
and \cref{lem:fullcond} with (A1) turns the right-hand side into the ratio $R_{ij}$ of four TFM predictives that the
algorithm computes with two forward passes. A swap preserves the counts, so $\lik(y')=\lik(y)$ and $R_{ij}=\pi(y')/\pi(y)$,
independently of the order of $(i,j)$. The number of differently labelled pairs depends only on the counts, so the
proposal is symmetric and the chain is Metropolis with stationary distribution $\pi$. Two assignments of a bag with
the same counts are joined by a sequence of swaps preserving the counts, each with positive probability by
\cref{cor:positivity}; bags are independent, so the product chain is irreducible (the lazy version is also aperiodic).
\end{proof}

\begin{proof}[Proof of \cref{thm:exact}(c)]
(Inductive.) The ergodic theorem for a finite irreducible chain applied to the bounded function
$y\mapsto p(y_*\mid x_*,(X,y))$, with (A1), gives almost sure convergence of \eqref{eq:predict} to $\bar p$.
By the law of total probability and (A0), $p(y_*\mid x_*,X,\W)=\sum_yp(y\mid x_*,X,\W)\,p(y_*\mid x_*,(X,y))$; the
difference of the expectations of one $[0,1]$-valued function under two measures is at most their TV distance.
If the prior is separable, $\phi=(\phi_x,\phi_{y\mid x})$ independent a priori with $p(x\mid\phi)=p(x\mid\phi_x)$ and
$p(y\mid x,\phi)=p(y\mid x,\phi_{y\mid x})$, then $p(y\mid x_*,X,\W)\propto\lik(y)\int p(y\mid X,\phi_{y\mid x})p(\phi_{y\mid x})d\phi_{y\mid x}$
times a factor that does not depend on $y$, so $p(y\mid x_*,X,\W)=p(y\mid X,\W)$ and the bound is zero.
(Transductive.) Adding $x_*$ as a free point with $\lik_*\equiv1$ gives a chain for $p(y,y_*\mid X,x_*,\W)$ that is exact by
part (a) (for LLP, the swap chain combined with single-site updates of $x_*$, each preserving the target). By
\cref{lem:fullcond}, the Rao--Blackwell term is $\tfm(y_*\mid x_*,(X,y^{(s)}))$, whose average converges to
$p(y_*\mid x_*,X,\W)$.
\end{proof}

\begin{proof}[Proof of \cref{prop:bayes}]
By the tower property, $\E[S(f,y_*)]=\E\big[\E[S(f,y_*)\mid X,\W,x_*]\big]$. For each value of $(X,\W,x_*)$, $f$ is a fixed
distribution and $y_*\sim p(\cdot\mid x_*,X,\W)$; strict propriety gives the unique minimiser
$f=p(\cdot\mid x_*,X,\W)$. Its risk is finite (at most $\log C$ for the log score).
\end{proof}

\subsection{Theorem 2}
\begin{proof}[Proof of \cref{thm:stability}]
\emph{Existence and uniqueness.} The state space is finite, so $\tilde P$ has a stationary distribution. With $\tfm>0$
every single-site update and every MH acceptance probability is positive, so $\tilde P$ has the transition graph of $P$
and is irreducible; $\tilde\pi$ is unique. Compatibility of the conditionals of $\tfm$ is never used.
\emph{Perturbation identity.} Let $\Delta=\tilde\pi-\pi$ and $\nu=\tilde\pi(\tilde P-P)$, so $\nu(\Omega)=0$. From
$\tilde\pi=\tilde\pi\tilde P$ and $\pi=\pi P$, $\Delta=\nu+\Delta P$, and iterating, $\Delta=\sum_{t<T}\nu P^t+\Delta P^T$.
\emph{Contraction.} For a signed measure $\mu$ with $\mu(\Omega)=0$, $\lVert\mu K\rVert_1\le\delta(K)\lVert\mu\rVert_1$.
If $\tau(P)<\infty$ then $\delta(P^T)\to0$, the remainder vanishes, and $\lVert\Delta\rVert_1\le\lVert\nu\rVert_1\tau(P)$; if
$\tau(P)=\infty$ there is nothing to prove (no aperiodicity is needed). Finally
$\lVert\nu\rVert_1\le\sum_y\tilde\pi(y)\lVert\tilde P(y,\cdot)-P(y,\cdot)\rVert_1=2\bar\varepsilon$.
\emph{Single-site error.} $\tilde P(y,\cdot)-P(y,\cdot)$ averages the differences of the single-site updates, and each is the
TV distance between $\tfm\lik_i$ and $\ppd\lik_i$ after normalisation. With $Z_p=\sum_k\ppd(k)\lik_i(k)$,
$2\,\TV(\tfm_\lik,\ppd_\lik)\le\big(\sum_k\lik_i(k)\lvert\tfm(k)-\ppd(k)\rvert+\lvert Z_q-Z_p\rvert\big)/Z_p\le4(\max_k\lik_i(k)/Z_p)\,\TV(\tfm,\ppd)$,
so the error is at most $2\kappa\varepsilon_{\TV}$ with $\kappa=\sup\max_k\lik_i(k)/Z_p$.
\emph{Prediction.} $\lvert\E_{\tilde\pi}\tfm-\E_\pi\ppd\rvert\le\lvert\E_{\tilde\pi}(\tfm-\ppd)\rvert+\lvert\E_{\tilde\pi}\ppd-\E_\pi\ppd\rvert\le\varepsilon'_{\TV}+\TV(\tilde\pi,\pi)$.
\end{proof}
\emph{Relation to classical bounds.} With $E=\tilde P-P$ and the fundamental matrix $Z$, $\tilde\pi-\pi=\nu Z$
\citep{schweitzer1968,meyer1980}, which gives $\TV(\tilde\pi,\pi)\le\tau_1(Z)\bar\varepsilon$ with $\tau_1(Z)\le\tau(P)$
\citep{seneta1988,cho2001}; \eqref{eq:stability} is the weaker, more interpretable form. It is attained by a two-state
chain. For a single-site sweep the per-sweep error is at most $\lvert F\rvert$ times the per-update error; the MH swap
kernel has error at most $4\varepsilon_\infty$, with $\varepsilon_\infty=\sup\lvert\log\tfm-\log\ppd\rvert$.

\subsection{Block updates}
\begin{proposition}[Blocks]\label{prop:blocks}
(a) A block update whose proposal is the product of block-wise TFM predictives (for LLP, conditioned on bag counts),
accepted with the MH ratio computed by Brook's lemma along the changed labels, leaves $\pi$ invariant when $\tfm=\ppd$.
(b) Without the MH step, \cref{thm:stability} applies with an error $\bar\varepsilon_B$ that is bounded only under an
extra stability assumption on the PPD, and the bound grows with $n$ for fixed $K$.
\end{proposition}
\begin{proof}[Proof of \cref{prop:blocks}]
(a) The proposal $Q_B$ depends only on $y_{-B}$, which the step does not change, so the update is MH-within-Gibbs with
an independent proposal. Brook's lemma along the changed labels with \cref{lem:fullcond} gives the exact ratio
$p(y'\mid X)/p(y\mid X)$, multiplied by the ratio of $\lik$ (one for LLP). The proposal is positive on the support of the
block and acceptance is positive, so the chain is irreducible.
(b) Let $P_{\mathrm{bG}}$ be exact block Gibbs. Under the stability assumption that adding one labelled point or one
unlabelled feature to a context changes the PPD by at most $\beta$ in TV, the chain rule and a hybrid argument give
$\TV\big(\prod_{i\in B\cap F}\ppd(\cdot\mid x_i,D_{-B}),\,p(y_{B\cap F}\mid y_{-(B\cap F)},X)\big)\le m_f(m-1)\beta$
for a block of $m$ points of which $m_f$ are free, and after reweighting by $\lik$,
$\bar\varepsilon_B\le2\kappa_B\,m_f(m-1)\beta$; \cref{thm:stability} with $P=P_{\mathrm{bG}}$ concludes.
\end{proof}
\emph{Why this is vacuous at our sizes.} With $\beta=O(1/n)$ and fixed $K$ ($m=n/K$), $m_f(m-1)\beta=O(n/K^2)$ grows
with $n$. For bag-blocks, $\kappa_B$ is the inverse probability that all bags of a group match their counts under the
product proposal: for $n\approx2{,}000$ and $K=10$, $\log_{10}\kappa_B\approx15.8$ for bags of 8. On an exactly solvable
LLP problem (3 bags of 4), the uncorrected bag-block chain has $\TV(\tilde\pi,\pi)$ between 0 and about 0.6 depending on
the prior and grouping, so the approximation is real at small $n$; at our sizes we measure its effect on predictions
instead (\cref{sec:ablation}).

\subsection{When the naive TFM suffices}
\paragraph{Complement to \cref{prop:naive}(a).} For binary labels with class prior $\varpi$, for any score $s$,
$\mathrm{AUC}(\eta)-\mathrm{AUC}(s)\le\E\lvert s-\tilde\eta\rvert/(\lambda\varpi(1-\varpi))$: the AUC regret of a learner fitted
to naive labels is its own estimation error, amplified by $1/\lambda$.
\begin{proof}[Proof of \cref{prop:naive} and of the complement]
(a) With $P(\tilde w=k\mid x,y)=M_{y,k}$ independent of $x$ and $M_{y,k}=b_k$ for $y\neq k$,
$\tilde\eta_k=M_{k,k}\eta_k+b_k(1-\eta_k)=b_k+\lambda_k\eta_k$, $\lambda_k=M_{k,k}-b_k$. For PLL with each wrong class
entering the candidate set independently with probability $q$ and the naive label chosen uniformly from the set,
$M_{y,y}=a=\E[1/(1+N)]=\frac{1-(1-q)^C}{Cq}$ with $N\sim\mathrm{Bin}(C-1,q)$ and $b=(1-a)/(C-1)$, so $\lambda=(Ca-1)/(C-1)>0$
for $q<1$. For LLP with random bags of size $k$, $P(\tilde w_i=1\mid x_i,y_i)=(y_i+(k-1)\varpi)/k$, so $\lambda=1/k$.
A strictly increasing affine map preserves ranks, hence one-vs-rest AUC. For the binary bound, write
$\varpi(1-\varpi)\mathrm{AUC}(s)=\E[\eta(X)(1-\eta(X'))h_s(X,X')]$ for independent copies, with
$h_s=\mathbb 1[s(X)>s(X')]+\frac12\mathbb 1[s(X)=s(X')]$; symmetrising,
$\varpi(1-\varpi)[\mathrm{AUC}(\eta)-\mathrm{AUC}(s)]=\frac12\E\big[\lvert\Delta\rvert(\mathbb 1[s\text{ reverses the pair}]+\frac12\mathbb 1[s\text{ ties},\Delta\neq0])\big]$
with $\Delta=\eta(X)-\eta(X')$. With $t=(s-b)/\lambda$, which has the ranks of $s$, a reversed or tied pair has
$\lvert\Delta\rvert\le\lvert\eta(X)-t(X)\rvert+\lvert\eta(X')-t(X')\rvert$, so the right-hand side is at most
$\E\lvert s-\tilde\eta\rvert/\lambda$. For $\lambda=1$ and scores that are class-probability estimates, this is Corollary~12 of \citet{agarwal2014}.
(b) $p_\theta(\tilde w)=\sum_Op_\theta(O)R(\tilde w\mid O)$ gives
$\partial_\theta\log p_\theta(\tilde w)=\E_\theta[\partial_\theta\log p_\theta(O)\mid\tilde w]$, and Jensen's inequality for
conditional expectations gives $I_{\tilde w}(\theta)\le I_O(\theta)$, with equality iff the score of $O$ is
$\sigma(\tilde w)$-measurable.
\end{proof}
\emph{Information lost by the naive labels} (logistic model $\eta=\sigma(\theta x)$, $\theta=2$, $x\sim\mathcal N(0,I)$; per
point). LLP: the Godambe information of the pointwise (composite) likelihood on naive labels, relative to the
information of the bag counts, is 0.134, 0.060, 0.033, 0.017, 0.009 for bags of 2, 4, 8, 16, 32. PLL: the naive labels
retain 0.24--0.26 of the information in the candidate sets at $q=0.3$ and 0.06--0.10 at $q=0.6$ ($C\in\{3,4,7\}$).
For PU-SCAR and label noise the ratio is 1.


\FloatBarrier
\section{Pre-registration and timeline}\label{app:prereg}
Both studies were run under a written protocol with hypotheses, comparators, tests, datasets and decision rules fixed
before any evaluation data were used; every later change was logged as a dated deviation. In the anonymised repository
each freeze is a commit, and every result file records the hashes of the code that produced it.

\begin{table}[!htbp]
\centering
\caption{Timeline of the two studies (2026, UTC). ``Exploratory'' marks analyses that were not part of the registered
tests.}
\label{tab:timeline}
\small
\begin{tabular}{@{}lp{0.72\textwidth}@{}}
\toprule
Date & Event \\
\midrule
26 Oct & Study~1 protocol frozen (hypotheses H1--H5, datasets, baselines, tuning spaces, primary-comparator rule, tests,
ablations, sensitivity analysis). Later changes were logged as dated deviations: all but three before the evaluation
runs; the three made during the analysis (dropping an external LLP benchmark whose data were unavailable, fixing the
reading of an ambiguous rule to its weakest form, labelling the TFM ablation exploratory) are reported. \\
28 Oct--1 Nov & Study~1 runs on the evaluation datasets (TabPFN~v2; registered TabICL~v2 ablation; baselines). \\
30 Oct & Study~1 analysis: H2 and H4 not met, H3 and H5 met; TFM ablation passes all tests (exploratory).
Decision to run a confirmatory replication changing only the TFM. \\
1 Nov & Study~2 dataset census and selection from metadata only; flags for the sensitivity analysis set by hand. \\
1 Nov 15:19 & \textbf{Study~2 frozen} (commit), before any Study~2 table was downloaded. \\
1--2 Nov & Data download and post-check rule applied (two datasets fail $n\ge500$ after de-duplication; first reserve
added, subsampled to 10{,}000 rows). Exploratory response-map runs on Study~1 datasets only. \\
2--3 Nov & Study~2 runs; confirmatory analysis with the pre-written script. One post-hoc fix of a floating-point tie in
the analysis (a tie counted as a loss); no conclusion changed. Failed cells of one baseline rerun once. \\
3 Nov & Exploratory runs on Study~1 datasets: $\widehat R$ and MH-corrected blocks with TabICL~v2. \\
\bottomrule
\end{tabular}
\end{table}

\FloatBarrier
\section{Study 1 with TabICL v2 (registered TFM ablation)}\label{app:s1-tabicl}
\Cref{tab:s1tabicl} applies the H2 and H4 tests of Study~1 to the registered TFM ablation, in which every TFM-based method
uses TabICL~v2. With clean labels the two TFMs are about equally accurate (oracle within \val{s1-tabicl-oracle-maxdiff}); the difference appears only
with weak labels. Because TabICL~v2 was promoted after these results were seen, they are exploratory; Study~2 is the
confirmatory test.
\begin{table}[!htbp]
\centering
\caption{Study~1 datasets, TabICL~v2 for all TFM-based methods (exploratory). Same tests and comparators as
\cref{tab:main}.}
\label{tab:s1tabicl}
\small
\resizebox{\textwidth}{!}{% generated by paper/scripts/make_appendix.py — do not edit
\begin{tabular}{@{}llrrrrrr@{}}
\toprule
Setting & Comparator & $\Delta$ [95\% CI] & W/T/L & $p_{\mathrm{Holm}}$ & \multicolumn{3}{c}{mean: \ours{} (TabICL / TabPFN) / oracle (TabICL)} \\
\midrule
LLP, bag 8 & DLLP (linear) & +0.058 [+0.029, +0.093] & 10/0/0 & 0.004 & 0.889 & 0.855 & 0.928 \\
PLL, $q=0.3$ & CC, tuned & +0.029 [+0.013, +0.051] & 12/0/2 & 0.004 & 0.854 & 0.847 & 0.865 \\
PLL, $q=0.6$ & CC, tuned & +0.038 [+0.013, +0.071] & 12/0/2 & 0.009 & 0.831 & 0.661 & 0.865 \\
\bottomrule
\end{tabular}
}
\end{table}

\FloatBarrier
\section{Study 2: sensitivity analysis, all baselines and per-dataset results}\label{app:s2-s1}
\paragraph{Sensitivity.} Five datasets were flagged before any data were downloaded (repeated measurements of the same
units, or temporal drift): SpeedDating and eye\_movements (binary); JapaneseVowels, cjs and gas-drift (multiclass).
Without them, every conclusion is unchanged (\cref{tab:s2s1}). With seven multiclass datasets, Holm-adjusted Wilcoxon
tests cannot reach $0.05$ unless every comparison is won; the registered rule for this case checks that the mean
difference keeps its sign and that \ours{} still wins on at least 60\% of the remaining datasets, which holds at both
levels.
\begin{table}[!htbp]
\centering
\caption{Study~2 without the flagged datasets.}
\label{tab:s2s1}
\small
\resizebox{\textwidth}{!}{% generated by paper/scripts/make_appendix.py — do not edit
\begin{tabular}{@{}llrrrrrr@{}}
\toprule
Setting & Comparator & datasets & $\Delta$ (all datasets) & $\Delta$ (flagged removed) & W/T/L & $p_{\mathrm{Holm}}$ & $\Delta$ vs naive \\
\midrule
LLP, bag 8 & DLLP (linear) & 16 & +0.052 & +0.051 & 15/0/1 & \ensuremath{<}0.001 & +0.124 \\
PLL, $q=0.3$ & CC, tuned & 7 & +0.012 & +0.016 & 6/1/0 & 0.062 & +0.023 \\
PLL, $q=0.6$ & CC, tuned & 7 & +0.012 & +0.013 & 5/1/1 & 0.062 & +0.168 \\
\bottomrule
\end{tabular}
}
\end{table}

\paragraph{All baselines.} Against each baseline separately (Holm over all baseline--level pairs of a setting),
\ours{} is significantly better than all ten LLP baselines; on PLL every mean difference is positive and 7 of 16 pairs are
significant after this stricter correction (\cref{tab:s2perbase}). One baseline (default PRODEN) produced non-finite
probabilities on 5 of 6 gas-drift cells, also on a rerun; its means use the remaining cells.
\begin{table}[!htbp]
\centering
\caption{Study~2, \ours{} against every baseline. Bold: significant in favour of \ours{}.}
\label{tab:s2perbase}
\small
% generated by paper/scripts/make_appendix.py — do not edit
\begin{tabular}{@{}llrrrr@{}}
\toprule
Setting & Baseline & datasets & $\Delta$ & W/T/L & $p_{\mathrm{Holm}}$ \\
\midrule
LLP, bag 8 (AUC) & DLLP (MLP) & 18 & +0.057 & 17/0/1 & \textbf{\ensuremath{<}0.001} \\
 & DLLP (MLP), tuned & 18 & +0.041 & 16/0/2 & \textbf{\ensuremath{<}0.001} \\
 & DLLP (linear) & 18 & +0.052 & 16/0/2 & \textbf{\ensuremath{<}0.001} \\
 & DLLP (linear), tuned & 18 & +0.046 & 16/0/2 & \textbf{\ensuremath{<}0.001} \\
 & LLP-VAT & 18 & +0.062 & 16/0/2 & \textbf{\ensuremath{<}0.001} \\
 & LLP-VAT, tuned & 18 & +0.042 & 16/0/2 & \textbf{\ensuremath{<}0.001} \\
 & EasyLLP & 18 & +0.116 & 18/0/0 & \textbf{\ensuremath{<}0.001} \\
 & EasyLLP, tuned & 18 & +0.069 & 17/0/1 & \textbf{\ensuremath{<}0.001} \\
 & GLWS & 18 & +0.056 & 17/1/0 & \textbf{\ensuremath{<}0.001} \\
 & GLWS, tuned & 18 & +0.051 & 18/0/0 & \textbf{\ensuremath{<}0.001} \\
\addlinespace
PLL, $q=0.3$ (acc.) & PRODEN & 9 & +0.048 & 9/0/0 & \textbf{0.043} \\
 & PRODEN, tuned & 10 & +0.027 & 8/1/1 & 0.070 \\
 & GLWS/ILL & 10 & +0.048 & 10/0/0 & \textbf{0.031} \\
 & GLWS/ILL, tuned & 10 & +0.014 & 10/0/0 & \textbf{0.031} \\
 & RC & 10 & +0.034 & 9/0/1 & \textbf{0.043} \\
 & RC, tuned & 10 & +0.013 & 8/1/1 & 0.070 \\
 & CC & 10 & +0.035 & 10/0/0 & \textbf{0.031} \\
 & CC, tuned & 10 & +0.012 & 8/2/0 & 0.070 \\
\addlinespace
PLL, $q=0.6$ (acc.) & PRODEN & 10 & +0.173 & 10/0/0 & \textbf{0.031} \\
 & PRODEN, tuned & 10 & +0.031 & 8/0/2 & 0.070 \\
 & GLWS/ILL & 10 & +0.128 & 10/0/0 & \textbf{0.031} \\
 & GLWS/ILL, tuned & 10 & +0.015 & 8/0/2 & 0.070 \\
 & RC & 10 & +0.048 & 8/0/2 & 0.070 \\
 & RC, tuned & 10 & +0.012 & 8/1/1 & 0.070 \\
 & CC & 10 & +0.047 & 8/0/2 & 0.070 \\
 & CC, tuned & 10 & +0.012 & 8/1/1 & 0.070 \\
\bottomrule
\end{tabular}

\end{table}

\begin{table}[!htbp]
\centering
\caption{Study~2, LLP with bags of 8: AUC per dataset (mean over 3 seeds). Bold: best non-oracle method among those
shown.}
\label{tab:s2llpds}
\small
\resizebox{\textwidth}{!}{% generated by paper/scripts/make_appendix.py — do not edit
\begin{tabular}{@{}lrrrrrrr@{}}
\toprule
Dataset & Gibbs & hard-EM & naive & DLLP-lin & DLLP tuned & LLP-VAT tuned & oracle \\
\midrule
ada & \textbf{0.898} & 0.895 & 0.742 & 0.888 & 0.880 & 0.883 & 0.912 \\
Bank\_Customer\_Churn & 0.849 & \textbf{0.855} & 0.785 & 0.750 & 0.837 & 0.834 & 0.877 \\
compas-two-years & \textbf{0.683} & 0.671 & 0.539 & 0.677 & 0.650 & 0.662 & 0.692 \\
eye\_movements & 0.564 & \textbf{0.589} & 0.538 & 0.569 & 0.549 & 0.548 & 0.701 \\
Fitness\_Club & \textbf{0.831} & 0.784 & 0.609 & 0.743 & 0.719 & 0.753 & 0.831 \\
hazelnut-contaminant & \textbf{0.993} & 0.989 & 0.853 & 0.941 & 0.953 & 0.951 & 0.995 \\
heloc & \textbf{0.797} & 0.788 & 0.759 & 0.789 & 0.788 & 0.787 & 0.813 \\
Is-this-a-good-customer & 0.633 & 0.634 & 0.576 & 0.610 & \textbf{0.656} & 0.631 & 0.737 \\
jasmine & \textbf{0.870} & 0.841 & 0.780 & 0.779 & 0.807 & 0.792 & 0.902 \\
Marketing\_Campaign & \textbf{0.878} & 0.870 & 0.667 & 0.609 & 0.630 & 0.635 & 0.915 \\
mushroom & \textbf{1.000} & \textbf{1.000} & 0.721 & 1.000 & 0.996 & 1.000 & 1.000 \\
NATICUSdroid & 0.983 & \textbf{0.985} & 0.911 & 0.973 & 0.976 & 0.977 & 0.988 \\
philippine & 0.774 & \textbf{0.781} & 0.711 & 0.681 & 0.699 & 0.700 & 0.934 \\
profb & 0.580 & 0.495 & 0.507 & 0.578 & 0.569 & \textbf{0.587} & 0.675 \\
scene & 0.971 & \textbf{0.983} & 0.823 & 0.973 & 0.977 & 0.973 & 0.991 \\
socmob & \textbf{0.980} & 0.979 & 0.794 & 0.956 & 0.950 & 0.947 & 0.987 \\
SpeedDating & \textbf{0.842} & 0.752 & 0.768 & 0.708 & 0.783 & 0.745 & 0.871 \\
sylvine & 0.993 & \textbf{0.996} & 0.955 & 0.960 & 0.965 & 0.958 & 0.997 \\
\midrule
mean & 0.840 & 0.827 & 0.724 & 0.788 & 0.799 & 0.798 & 0.879 \\
\bottomrule
\end{tabular}
}
\end{table}

\begin{table}[!htbp]
\centering
\caption{Study~2, PLL: accuracy per dataset (mean over 3 seeds). Bold: best non-oracle method among those shown.}
\label{tab:s2pllds}
\small
\resizebox{\textwidth}{!}{% generated by paper/scripts/make_appendix.py — do not edit
\begin{tabular}{@{}lrrrrrrrrrrrr@{}}
\toprule
 & \multicolumn{6}{c}{$q=0.3$} & \multicolumn{6}{c}{$q=0.6$} \\
\cmidrule(lr){2-7}\cmidrule(l){8-13}
Dataset & Gibbs & hard-EM & naive & CC tuned & RC tuned & oracle & Gibbs & hard-EM & naive & CC tuned & RC tuned & oracle \\
\midrule
anneal & \textbf{0.991} & 0.989 & 0.955 & 0.976 & 0.972 & 0.996 & \textbf{0.977} & 0.945 & 0.801 & 0.957 & 0.962 & 0.996 \\
cardiotocography & \textbf{1.000} & \textbf{1.000} & 0.994 & \textbf{1.000} & \textbf{1.000} & 1.000 & \textbf{1.000} & 0.886 & 0.476 & \textbf{1.000} & \textbf{1.000} & 1.000 \\
cjs & 0.997 & 0.997 & 0.962 & 0.997 & \textbf{0.998} & 1.000 & \textbf{0.993} & 0.830 & 0.391 & 0.991 & 0.989 & 1.000 \\
gas-drift & 0.996 & \textbf{0.997} & 0.993 & 0.991 & 0.992 & 0.998 & \textbf{0.995} & 0.994 & 0.930 & 0.988 & 0.986 & 0.998 \\
JapaneseVowels & 0.991 & \textbf{0.995} & 0.968 & 0.984 & 0.980 & 0.998 & \textbf{0.989} & 0.983 & 0.785 & 0.970 & 0.973 & 0.998 \\
MIC & \textbf{0.887} & 0.886 & 0.883 & 0.873 & 0.872 & 0.891 & \textbf{0.886} & 0.877 & 0.848 & 0.870 & 0.871 & 0.891 \\
students-dropout & 0.777 & \textbf{0.782} & 0.775 & 0.754 & 0.754 & 0.786 & \textbf{0.770} & 0.764 & 0.744 & 0.734 & 0.733 & 0.786 \\
website\_phishing & \textbf{0.823} & 0.820 & 0.792 & 0.787 & 0.787 & 0.844 & 0.786 & \textbf{0.787} & 0.735 & 0.778 & 0.778 & 0.844 \\
wine-quality-white & 0.563 & \textbf{0.571} & 0.538 & 0.547 & 0.546 & 0.585 & \textbf{0.556} & 0.510 & 0.481 & 0.539 & 0.537 & 0.585 \\
yeast & \textbf{0.597} & 0.578 & 0.543 & 0.588 & 0.591 & 0.606 & 0.561 & 0.456 & 0.275 & \textbf{0.563} & 0.563 & 0.606 \\
\midrule
mean & 0.862 & 0.862 & 0.840 & 0.850 & 0.849 & 0.870 & 0.851 & 0.803 & 0.647 & 0.839 & 0.839 & 0.870 \\
\bottomrule
\end{tabular}
}
\end{table}

\begin{table}[!htbp]
\centering
\caption{Study~2: mean score of every method (datasets averaged over 3 seeds). PRODEN (default) failed on gas-drift in
5 of 6 cells (non-finite probabilities, reproduced on a rerun); its $q=0.3$ mean is over 9 datasets and its $q=0.6$
gas-drift entry uses one seed.}
\label{tab:s2means}
\small
\resizebox{\textwidth}{!}{% generated by paper/scripts/make_tables.py — do not edit
\begin{tabular}{@{}lrr@{\hspace{2.2em}}lrrrr@{}}
\toprule
\multicolumn{3}{@{}l}{LLP, bag 8 (AUC, 18 datasets)} & \multicolumn{5}{l}{PLL (accuracy, 10 datasets)} \\
\cmidrule(r){1-3}\cmidrule(l){4-8}
Method & default & tuned & Method & \multicolumn{2}{c}{$q=0.3$} & \multicolumn{2}{c}{$q=0.6$} \\
 & & & & default & tuned & default & tuned \\
\midrule
\textbf{In-Context Gibbs (ours)} & 0.840 &  & \textbf{In-Context Gibbs (ours)} & 0.862 &  & 0.851 &  \\
Hard-EM (same TFM) & 0.827 &  & Hard-EM (same TFM) & 0.862 &  & 0.803 &  \\
Stochastic EM (same TFM) & 0.835 &  & Stochastic EM (same TFM) & 0.858 &  & 0.843 &  \\
Naive TFM & 0.724 &  & Naive TFM & 0.840 &  & 0.647 &  \\
DLLP (MLP) & 0.783 & 0.799 & PRODEN & 0.800 & 0.836 & 0.678 & 0.820 \\
DLLP (linear) & 0.788 & 0.794 & GLWS & 0.814 & 0.848 & 0.724 & 0.837 \\
LLP-VAT & 0.778 & 0.798 & RC & 0.828 & 0.849 & 0.803 & 0.839 \\
EasyLLP & 0.724 & 0.771 & CC & 0.827 & 0.850 & 0.804 & 0.839 \\
GLWS & 0.784 & 0.789 & Oracle (true labels) & 0.870 &  & 0.870 &  \\
Oracle (true labels) & 0.879 &  &  &  &  &  &  \\
\bottomrule
\end{tabular}
}
\end{table}

\FloatBarrier
\section{Correctness experiments}\label{app:rq1}
We enumerate every labelling of small problems ($n\le14$), build the exact transition matrix of each kernel with an
analytic PPD as $\tfm$, and compare its stationary distribution with $\pi$. Two priors are used: one in which features
carry no information about the label mechanism, and a three-component mixture in which they do. Across all
\val{rq1-e1-n} configurations the TV distance is at most $\val{rq1-e1-tvmax}$ (\cref{tab:rq1e1}), and running the
actual sampler for $10^5$ sweeps reaches the Monte Carlo floor of i.i.d.\ sampling. For \cref{thm:stability} we use
conditionals with deliberate, incompatible perturbations (\val{rq1-e3-n} configurations) and small PFNs trained on the same
prior; the bound held in \val{rq1-thm2-held} of \val{rq1-thm2-n} configurations and was \val{rq1-loose-min} to
\val{rq1-loose-max} times loose (median \val{rq1-loose-median}; \cref{tab:rq1thm2}); a fully trained PFN gives TV
\val{rq1-pfn-final-single} (single-site) and \val{rq1-pfn-final-swap} (swap). On an exactly solvable pooled problem with
$n=\val{rq1-e5-nmax}$ and $K=10$, blocking over all points changes predictions by $\val{rq1-e5-pred-all-n1000}$, whereas
blocking only the free points biases them by \val{rq1-e5-pred-free-n1000}, which is why blocks include fixed points.
\begin{table}[!htbp]
\centering
\caption{Exactness (\cref{thm:exact}, \cref{prop:blocks}(a)): TV between the stationary distribution of the exact kernel
and $\pi$.}
\label{tab:rq1e1}
\small
% generated by paper/scripts/make_appendix.py — do not edit
\begin{tabular}{@{}lllrrr@{}}
\toprule
Prior & Weak labels & Kernel & $n$ & states & $\TV(\text{stationary},\pi)$ \\
\midrule
cell & PU-SCAR (c=0.5) & single & 10 & 128 & 3.3e-15 \\
cell & SSL & single & 10 & 128 & 4.0e-15 \\
cell & PLL & single & 10 & 128 & 4.0e-15 \\
cell & Noisy (T known) & single & 8 & 256 & 1.9e-14 \\
cell & LLP (3 bags of 4) & swap & 12 & 96 & 3.1e-15 \\
cell & SSL, block-MH (m=3) & block-mh & 10 & 128 & 2.8e-14 \\
mixture & PU-SCAR (c=0.5) & single & 10 & 128 & 3.3e-15 \\
mixture & SSL & single & 10 & 128 & 3.0e-15 \\
mixture & PLL & single & 10 & 128 & 3.0e-15 \\
mixture & Noisy (T known) & single & 8 & 256 & 2.3e-14 \\
mixture & LLP (3 bags of 4) & swap & 12 & 96 & 4.8e-15 \\
mixture & SSL, block-MH (m=3) & block-mh & 10 & 128 & 4.0e-15 \\
cell & PLL (C=3, $|$S\_i$|$=2) & single & 7 & 128 & 1.2e-14 \\
mixture & PLL (C=3, $|$S\_i$|$=2) & single & 7 & 128 & 3.8e-15 \\
\bottomrule
\end{tabular}

\end{table}
\begin{table}[!htbp]
\centering
\caption{Stability (\cref{thm:stability}): actual TV between the stationary distributions of the approximate and exact
chains, the bound $\tau(P)\bar\varepsilon$, and their ratio (looseness).}
\label{tab:rq1thm2}
\small
\resizebox{\textwidth}{!}{% generated by paper/scripts/make_appendix.py — do not edit
\begin{tabular}{@{}llrrrrr@{}}
\toprule
Conditionals & Kernel & configs & median $\TV(\tilde\pi,\pi)$ & median bound & bound held & median looseness \\
\midrule
perturbed PPD (incompatible) & single & 15 & 0.0181 & 0.246 & 15/15 & 9.8 \\
perturbed PPD (incompatible) & swap & 15 & 0.0254 & 0.066 & 15/15 & 2.2 \\
small PFN (trained) & single & 6 & 0.1419 & 0.628 & 6/6 & 4.9 \\
small PFN (trained) & swap & 6 & 0.0480 & 0.107 & 6/6 & 3.3 \\
\bottomrule
\end{tabular}
}
\end{table}

\FloatBarrier
\section{Ablations}\label{app:ablations}
\begin{table}[!htbp]
\centering
\caption{Mixing and the exact sampler, by TFM. TabPFN: Study~1 (registered ablations); TabICL: exploratory lot on the
Study~1 datasets. ``3 small'': the three smallest datasets of each group, as registered for $\widehat R$. Columns 2--5:
medians over cells of the fraction of training labels with $\widehat R>1.1$ (two chains) and of the absolute difference
between the test metrics of the two chains. Columns 6--7: mean difference in AUC (LLP) or accuracy (PLL) between the
MH-corrected and the default sampler at 30 sweeps, with the mean MH acceptance rate. In the last row, $\widehat R$ uses all
four datasets with $\geq$8 classes, and the MH comparison only two of them (mfeat-factors, optdigits; 5 of 12 cells),
because the corrected sampler ran out of the 12-hour compute session on optdigits.}
\label{tab:s2e}
\small
\resizebox{\textwidth}{!}{% generated by paper/scripts/make_tables.py — do not edit
\begin{tabular}{@{}lrrrrrr@{}}
\toprule
 & \multicolumn{2}{c}{labels with $\widehat R>1.1$} & \multicolumn{2}{c}{test-metric gap, 2 chains} & \multicolumn{2}{c}{MH-corrected $-$ default} \\
\cmidrule(lr){2-3}\cmidrule(lr){4-5}\cmidrule(l){6-7}
Setting (datasets) & TabPFN & TabICL & TabPFN & TabICL & TabPFN & TabICL [acc.\ rate] \\
\midrule
LLP, bag 8 (3 small) & 11\% & 28\% & 0.041 & 0.018 & +0.010 & +0.033 [0.68] \\
PLL, $q=0.3$ (3 small) & 6\% & 6\% & 0.006 & 0.018 & +0.006 & \ensuremath{-}0.009 [0.78] \\
PLL, $q=0.6$ (3 small) & 21\% & 31\% & 0.016 & 0.021 & -- & \ensuremath{-}0.021 [0.72] \\
PLL, $q=0.6$, $\geq$8 classes (4; MH: 2) & -- & 22\% & -- & 0.002 & -- & \ensuremath{-}0.016 [0.47] \\
\bottomrule
\end{tabular}
}
\end{table}

\begin{table}[!htbp]
\centering
\caption{\ours{} versus hard-EM with the same TFM and warm start, log loss. A cell is a (dataset, weak-label level)
pair; Bayesian signed-rank test with ROPE $\pm0.005$; Spearman correlation between weakness level and the gap.
$^\dagger$Computed on PLL cells (LLP has a single level), as stated before the analysis was run.}
\label{tab:h3}
\small
% generated by paper/scripts/make_tables.py — do not edit
\begin{tabular}{@{}lrrrrr@{}}
\toprule
Study (TFM; settings) & cells & Gibbs better & $P(\text{Gibbs better})$ & $P(\text{ROPE})$ & Spearman $\rho$ \\
\midrule
Study 1 (TabPFN v2; all six settings) & 128 & 86.7\% & 1.00 & 0.00 & 0.36 \\
Study 2 (TabICL v2; LLP bag 8, PLL) & 38 & 78.9\% & 1.00 & 0.00 & 0.66$^\dagger$ \\
\bottomrule
\end{tabular}

\end{table}

\begin{table}[!htbp]
\centering
\caption{Study~1 ablations (TabPFN~v2): each variant against the default on the same cells. Blocks, sweeps and MH: the
three smallest datasets per group; transductive and cold start: all datasets.}
\label{tab:ablation}
\small
% generated by paper/scripts/make_appendix.py — do not edit
\begin{tabular}{@{}llrrrr@{}}
\toprule
Variant (vs.\ default) & Setting & cells & $\Delta$ AUC/acc. & W/T/L & $\Delta$ log loss (+ = better) \\
\midrule
$K=20$ & LLP, bag 8 & 9 & +0.007 & 5/0/4 & +0.024 \\
$K=20$ & PLL, 0.3 & 9 & +0.006 & 5/2/2 & \ensuremath{-}0.000 \\
$K=5$ (default 10) & LLP, bag 8 & 9 & \ensuremath{-}0.005 & 6/0/3 & +0.006 \\
$K=5$ (default 10) & PLL, 0.3 & 9 & +0.004 & 5/1/3 & +0.009 \\
15 sweeps (default 30) & LLP, bag 8 & 9 & +0.005 & 4/0/5 & +0.001 \\
15 sweeps (default 30) & PLL, 0.3 & 9 & +0.004 & 4/2/3 & \ensuremath{-}0.001 \\
60 sweeps & LLP, bag 8 & 9 & +0.014 & 4/0/5 & \ensuremath{-}0.012 \\
60 sweeps & PLL, 0.3 & 9 & +0.005 & 4/5/0 & +0.009 \\
MH-corrected blocks (exact) & LLP, bag 8 & 9 & +0.010 & 4/0/5 & \ensuremath{-}0.019 \\
MH-corrected blocks (exact) & PLL, 0.3 & 9 & +0.006 & 4/3/2 & +0.009 \\
transductive prediction & LLP, bag 8 & 30 & \ensuremath{-}0.013 & 7/0/23 & \ensuremath{-}0.042 \\
cold start & LLP, bag 8 & 30 & \ensuremath{-}0.005 & 10/0/20 & \ensuremath{-}0.002 \\
cold start & PLL, 0.3 & 42 & +0.000 & 11/23/8 & +0.000 \\
cold start & PLL, 0.6 & 42 & +0.012 & 14/9/19 & +0.023 \\
PU-SAR prior $N(0,1)$ & PU-SAR, 0.1 & 30 & \ensuremath{-}0.002 & 15/3/12 & \ensuremath{-}0.015 \\
PU-SAR prior $N(0,1)$ & PU-SAR, 0.3 & 30 & +0.001 & 10/1/19 & +0.002 \\
PU-SAR prior $N(0,4^2)$ & PU-SAR, 0.1 & 30 & \ensuremath{-}0.004 & 7/9/14 & +0.000 \\
PU-SAR prior $N(0,4^2)$ & PU-SAR, 0.3 & 30 & \ensuremath{-}0.001 & 10/12/8 & \ensuremath{-}0.016 \\
\bottomrule
\end{tabular}

\end{table}

\FloatBarrier
\section{Real partial labels}\label{app:real}
\begin{table}[!htbp]
\centering
\caption{Real partial labels (Study~1, TabICL~v2 for all TFM-based methods): accuracy averaged over 10 random 90/10
splits, and mean rank over all splits of the three datasets among the ten ranked methods (\ours{}, naive TFM and eight
baselines). Bold: best non-oracle accuracy.}
\label{tab:real}
\small
% generated by paper/scripts/make_tables.py — do not edit
\begin{tabular}{@{}lrrrr@{}}
\toprule
Method & BirdSong & Lost & MSRCv2 & mean rank \\
\midrule
\textbf{In-Context Gibbs (ours)} & \textbf{0.759} & \textbf{0.750} & \textbf{0.628} & 2.15 \\
Hard-EM (same TFM) & 0.730 & 0.665 & 0.627 & -- \\
Stochastic EM (same TFM) & 0.744 & 0.716 & 0.622 & -- \\
Naive TFM & 0.689 & 0.481 & 0.491 & 7.27 \\
\midrule
PRODEN (default) & 0.667 & 0.373 & 0.344 & 9.88 \\
PRODEN (tuned) & 0.707 & 0.748 & 0.465 & 5.12 \\
GLWS (default) & 0.722 & 0.712 & 0.582 & 3.98 \\
GLWS (tuned) & 0.731 & 0.747 & 0.490 & 4.03 \\
RC (default) & 0.719 & 0.720 & 0.424 & 6.00 \\
RC (tuned) & 0.702 & 0.721 & 0.478 & 5.23 \\
CC (default) & 0.719 & 0.717 & 0.425 & 6.02 \\
CC (tuned) & 0.700 & 0.723 & 0.479 & 5.32 \\
\midrule
Oracle (true labels) & 0.835 & 0.889 & 0.824 & -- \\
\bottomrule
\end{tabular}

\end{table}

\FloatBarrier
\section{Datasets, baselines and tuning}\label{app:data}
\paragraph{Datasets.} Study~1 uses every OpenML-CC18 dataset with real (non-synthetic) data, 500 to 10{,}000 rows, at
most 500 features and, for binary tasks, a minority class of at least 10\% (multiclass: 3--10 classes), one per source;
Study~2 applies the same criteria to five curated suites after excluding OpenML-CC18 and Study~1
(\cref{tab:datasets}). Duplicate rows are removed before splitting. TFMs receive ordinally encoded features; MLP
baselines receive median-imputed, standardised features with one-hot encoded categories, fitted on the training split.

\paragraph{Baselines and tuning.} Each baseline runs in the configuration of its original code or paper and tuned by 20
random configurations (the default among them) from the spaces below, scored by the method's own loss on a 20\%
validation split with weak labels only (split by bag for LLP); the best configuration is refit on the full training
split. MLP-based learners (DLLP, LLP-VAT, EasyLLP, GLWS, PRODEN) search learning rate $\sim\mathrm{LogU}(10^{-4},10^{-2})$,
weight decay $\sim\mathrm{LogU}(10^{-6},10^{-2})$, hidden layers in $\{(),(64),(128,128)\}$ and epochs in
$\{100,200,400\}$, plus bags per batch for LLP, VAT parameters for LLP-VAT, the optimiser for PRODEN and the entropy
term for GLWS/ILL; RC and CC use a linear model with the PRODEN space (\cref{tab:tuning}).

\begin{table}[!htbp]
\centering
\caption{Search spaces (20 random configurations, the default included). LogU: log-uniform.}
\label{tab:tuning}
\small
\resizebox{\textwidth}{!}{%
\begin{tabular}{@{}lll@{}}
\toprule
Baseline & Default & Search space \\
\midrule
DLLP (MLP, linear), EasyLLP & lr $10^{-3}$, wd 0, 200 epochs & lr $\sim$ LogU($10^{-4},10^{-2}$); wd $\sim$ LogU($10^{-6},10^{-2}$); hidden $\in\{(),(64),(128,128)\}$; epochs $\in\{100,200,400\}$; bags/batch $\in\{8,32\}$ \\
LLP-VAT & as DLLP + VAT defaults of the official code & as DLLP; $\xi\in\{10^{-6},10^{-3}\}$; $\epsilon\in\{0.1,0.5,1,2\}$; consistency weight $\sim$ LogU($10^{-2},1$) \\
GLWS / ILL & official code, MLP backbone & lr, wd as DLLP; hidden $\in\{(64),(128),(128,128)\}$; epochs $\in\{100,200,400\}$; entropy term on/off (PLL) \\
PRODEN & lr $10^{-3}$, wd $10^{-5}$, hidden $(128,128)$, 200 epochs, SGD & lr, wd, hidden, epochs as DLLP; optimiser $\in\{$SGD (momentum 0.9), Adam$\}$ \\
RC, CC & linear model as in the original paper & as PRODEN \\
SAR-EM (PU-SAR) & $C_f=C_e=1$ & $C_f,C_e\sim$ LogU($10^{-3},10^{3}$) \\
nnPU, uPU, forward correction & lr $10^{-3}$, wd 0, hidden $(100)$, 200 epochs & lr $\sim$ LogU($10^{-4},10^{-2}$); wd $\sim$ LogU($10^{-6},10^{-2}$); hidden $\in\{(64),(128),(128,128)\}$; epochs $\in\{100,200,400\}$ \\
VPU & official code & lr $\sim$ LogU($10^{-5},10^{-3}$); mixup $\lambda\in\{0.01,0.03,0.1,0.3\}$; $\alpha\in\{0.1,0.3,1\}$ \\
Self-training & logistic regression, $C=1$, threshold 0.75 & $C\sim$ LogU($10^{-3},10^{3}$); threshold $\in\{0.6,0.75,0.9,0.95\}$ \\
Label spreading & kNN, $k=7$, $\alpha=0.2$ & kernel $\in\{$kNN, RBF$\}$; $k\in\{3,5,7,11,15,21\}$; $\gamma\sim$ LogU($10^{-3},10$); $\alpha\in\{0.1,0.2,0.5,0.8\}$ \\
VIME & official code & $p_m\in\{0.1,0.3,0.5\}$; $\alpha\in\{0.5,1,2\}$; $\beta\in\{0.1,1,10\}$; hidden $\in\{(64),(128)\}$ \\
Confident learning + TFM & default (not tuned) & -- \\
\bottomrule
\end{tabular}}
\end{table}

\begin{table}[!htbp]
\centering
\caption{Datasets ($n$ after removing duplicates; $^\ast$subsampled to 10{,}000 rows by the registered rule; S1:
flagged for the sensitivity analysis).}
\label{tab:datasets}
\scriptsize
\resizebox{\textwidth}{!}{% generated by paper/scripts/make_appendix.py — do not edit
\begin{tabular}{@{}rlrrrc@{\hspace{1.2em}}rlrrrc@{}}
\toprule
\multicolumn{12}{@{}l}{Study 1 (OpenML-CC18)} \\
\midrule
ID & name & $n$ & $d$ & $C$ & S1 & ID & name & $n$ & $d$ & $C$ & S1 \\
\midrule
3 & kr-vs-kp & 3,196 & 36 & 2 &  & 28 & optdigits & 5,620 & 64 & 10 &  \\
29 & credit-approval & 690 & 15 & 2 &  & 46 & splice & 3,006 & 60 & 3 &  \\
44 & spambase & 4,210 & 57 & 2 &  & 182 & satimage & 6,430 & 36 & 6 &  \\
50 & tic-tac-toe & 958 & 9 & 2 &  & 188 & eucalyptus & 736 & 19 & 5 &  \\
1049 & pc4 & 1,344 & 37 & 2 &  & 458 & analcatdata\_authorship & 841 & 70 & 4 &  \\
1489 & phoneme & 5,395 & 5 & 2 &  & 469 & analcatdata\_dmft & 556 & 4 & 6 &  \\
1494 & qsar-biodeg & 1,052 & 41 & 2 &  & 1475 & first-order-theorem-proving & 5,636 & 51 & 6 &  \\
6332 & cylinder-bands & 540 & 37 & 2 &  & 1497 & wall-robot-navigation & 5,456 & 24 & 4 & \checkmark \\
23381 & dresses-sales & 500 & 12 & 2 &  & 1501 & semeion & 1,593 & 256 & 10 &  \\
38 & sick & 3,711 & 29 & 2 &  & 40966 & MiceProtein & 1,080 & 77 & 8 & \checkmark \\
11 & balance-scale & 625 & 4 & 3 & \checkmark & 40975 & car & 1,728 & 6 & 4 & \checkmark \\
12 & mfeat-factors & 1,994 & 216 & 10 &  & 40982 & steel-plates-fault & 1,941 & 27 & 7 &  \\
\bottomrule
\end{tabular}
\par\medskip
\begin{tabular}{@{}rlrrrc@{\hspace{1.2em}}rlrrrc@{}}
\toprule
\multicolumn{12}{@{}l}{Study 2 (new curated suites)} \\
\midrule
ID & name & $n$ & $d$ & $C$ & S1 & ID & name & $n$ & $d$ & $C$ & S1 \\
\midrule
24 & mushroom & 8,124 & 22 & 2 &  & 46930 & hazelnut-spread-contaminant-detection & 2,400 & 30 & 2 &  \\
312 & scene & 2,400 & 299 & 2 &  & 46938 & Is-this-a-good-customer & 1,723 & 13 & 2 &  \\
470 & profb & 672 & 9 & 2 &  & 46940 & Marketing\_Campaign & 2,058 & 25 & 2 &  \\
934 & socmob & 1,156 & 5 & 2 &  & 46969 & NATICUSdroid & 7,491 & 86 & 2 &  \\
40536 & SpeedDating & 8,378 & 120 & 2 & \checkmark & 181 & yeast & 1,453 & 8 & 10 &  \\
41143 & jasmine & 2,984 & 144 & 2 &  & 375 & JapaneseVowels & 9,961 & 14 & 9 & \checkmark \\
41145 & philippine & 5,832 & 308 & 2 &  & 1466 & cardiotocography & 2,115 & 35 & 10 &  \\
41146 & sylvine & 5,124 & 20 & 2 &  & 23380 & cjs & 1,811 & 33 & 6 & \checkmark \\
41156 & ada & 4,146 & 48 & 2 &  & 40498 & wine-quality-white & 3,961 & 11 & 7 &  \\
44130 & eye\_movements & 7,608 & 20 & 2 & \checkmark & 46906 & anneal & 886 & 38 & 5 &  \\
45026 & heloc & 9,444 & 22 & 2 &  & 46960 & students\_dropout\_and\_academic\_success & 4,424 & 36 & 3 &  \\
45039 & compas-two-years & 2,712 & 11 & 2 &  & 46963 & website\_phishing & 724 & 9 & 3 &  \\
46911 & Bank\_Customer\_Churn & 10,000 & 10 & 2 &  & 46980 & MIC & 1,699 & 111 & 8 &  \\
46927 & Fitness\_Club & 1,499 & 6 & 2 &  & 1476 & gas-drift & 10,000$^\ast$ & 128 & 6 & \checkmark \\
\bottomrule
\end{tabular}
\par\medskip
}
\end{table}

\FloatBarrier
\section{Secondary settings of Study 1}\label{app:secondary}
Study~1 also registered secondary settings (\cref{tab:secondary}). \ours{} improved strongly on the naive TFM for larger
bags, as for bags of 8, but not for PU with labels selected completely at random, class-dependent noise or SSL, where its
difference to the naive TFM was at most \val{s1-d3-max} AUC. This is the pattern \cref{prop:naive} predicts, and it was
registered in advance as exploratory analyses D1--D3: the advantage over the naive TFM grows from $q=0.3$ to $q=0.6$ (D1,
one-sided $p=\val{s1-d1-p}$) and from bags of 8 to 16 (D2, $p=\val{s1-d2-p}$), and stays below $+0.02$ AUC in the
settings where the naive labels lose no information (D3, all six levels). Against the primary comparators of these
settings no difference was significant except class-dependent noise at $r=0.4$, where confident learning combined with the
TFM was better.

\begin{table}[!htbp]
\centering
\caption{Study~1 secondary settings (TabPFN~v2). $\Delta$ positive favours \ours{}. $^\dagger$Log loss after matching the
known class prior (sign flipped so that positive is better).}
\label{tab:secondary}
\small
\resizebox{\textwidth}{!}{% generated by paper/scripts/make_appendix.py — do not edit
\begin{tabular}{@{}lllrrrrr@{}}
\toprule
Setting & Metric & Primary comparator & \multicolumn{3}{c}{\ours{} $-$ comparator} & \multicolumn{2}{c}{\ours{} $-$ naive TFM} \\
\cmidrule(lr){4-6}\cmidrule(l){7-8}
 & & & $\Delta$ & W/T/L & $p_{\mathrm{Holm}}$ & $\Delta$ & W/T/L \\
\midrule
LLP, bag 16 & AUC & DLLP (linear) & +0.040 & 7/0/3 & 0.232 & +0.206 & 9/0/1 \\
LLP, bag 32 & AUC & DLLP (linear) & +0.048 & 6/0/4 & 0.375 & +0.223 & 10/0/0 \\
PU-SAR, $\E[e\mid y=1]=0.1$ & AUC & nnPU & +0.034 & 6/0/4 & 0.967 & +0.007 & 6/0/4 \\
PU-SAR, $\E[e\mid y=1]=0.3$ & AUC & nnPU & +0.032 & 8/0/2 & 0.109 & +0.017 & 6/0/4 \\
PU-SCAR, $c=0.1$ & log loss$^\dagger$ & VPU, tuned & +0.118 & 8/0/2 & 0.111 & \ensuremath{-}0.033 & 5/0/5 \\
PU-SCAR, $c=0.3$ & log loss$^\dagger$ & VPU, tuned & +0.179 & 8/0/2 & 0.111 & \ensuremath{-}0.028 & 2/0/8 \\
Class-dependent noise, $r=0.2$ & AUC & CL + TFM & \ensuremath{-}0.006 & 3/0/7 & 0.105 & \ensuremath{-}0.013 & 1/0/9 \\
Class-dependent noise, $r=0.4$ & AUC & CL + TFM & \ensuremath{-}0.048 & 2/0/8 & 0.020 & \ensuremath{-}0.054 & 1/0/9 \\
SSL, 5\% labelled & AUC & Self-training, tuned & \ensuremath{-}0.002 & 5/0/5 & 1.000 & \ensuremath{-}0.041 & 0/0/10 \\
SSL, 20\% labelled & AUC & Self-training, tuned & \ensuremath{-}0.009 & 8/0/2 & 0.750 & \ensuremath{-}0.016 & 0/0/10 \\
\bottomrule
\end{tabular}
}
\end{table}

\FloatBarrier
\section{Statements}\label{app:statements}
\paragraph{Data and code availability.} All datasets are public (OpenML; the PALM collection for real partial labels).
Code, frozen protocols, raw per-cell results with code and data hashes, and the scripts that generate every number,
table and figure of this paper are available in an anonymised repository and will be released.

\paragraph{Compute.} All runs used free Kaggle sessions with two NVIDIA T4 GPUs (TFM-based methods) or CPUs (baselines).
Study~1 used about 42 hours of two-GPU sessions and 26 hours of four-process CPU sessions; Study~2 about 4 GPU-hours
and 80 CPU process-hours (two sessions of about ten hours with four processes each); the exploratory lots about 13 hours
of two-GPU sessions. Development and pilot runs before each freeze are not included.

\paragraph{Broader impacts.} Learning from weak labels lowers the cost of annotation and makes coarse or ambiguous labels
usable. Learning from label proportions also recovers individual-level predictions from data released only in
aggregate; when aggregation was chosen to protect individuals, such models can weaken that protection, and releases of
aggregated data should be assessed with this in mind. We have not measured how much individual-level information the
sampler recovers from bag counts. \todo{quantify, or keep this sentence (G9/REVIEW\_RESPONSE\_PLAN R8)}

\paragraph{Use of large language models.} An AI assistant was used for writing code, drafting proofs that the authors
checked, running analyses through the pre-written scripts, and drafting text. The authors designed the studies, approved
every pre-registration and decision, ran all experiments and are responsible for all content.
\todo{align with the venue's LLM policy; docs/AI\_USAGE.md}
\FloatBarrier
